(1) 均匀柔软不可伸长链条总长度 \(L\)、线密度 \(\mu\)，在无摩擦光滑圆角桌缘转向，一段长度 \(x(t)\) 竖直悬挂，其余部分在桌面上。忽略链条离开桌缘时的局部动力学损失。以 \(x\) 为广义坐标求动能与势能。
(2) 建立全链精确的一维运动方程，说明为何惯性质量不是 \(\mu x\)。
(3) 若从 \(x(0)=x_0>0\)、\(\dot x(0)=0\) 释放，求 \(x(t)\)、\(\dot x(t)\) 和到 \(x=L\) 的时间。
(4) 由能量法独立验证速度公式，并说明尖锐桌缘或耗散会导致哪些模型修正。
